% Original analytic teaching diagram, 2026-11-12.
% Covering the rationals of [0,1] listed as q_1=1/2, q_2=1/3, q_3=2/3, q_4=1/4, q_5=3/4, ...
% by intervals centred at q_k of length eps/2^k, eps = 0.2 (lengths 0.1, 0.05, 0.025, 0.0125, 0.00625).
% Total length at most eps, so the rationals have probability 0 under the uniform point.
\definecolor{VB}{HTML}{7188A3}\definecolor{VBF}{HTML}{E5EAF0}
\definecolor{VG}{HTML}{829C8B}\definecolor{VGF}{HTML}{E8EEE9}
\definecolor{VY}{HTML}{B4A16F}\definecolor{VYF}{HTML}{F1ECDD}
\definecolor{VR}{HTML}{A7787D}\definecolor{VRF}{HTML}{F0E5E6}
\definecolor{VP}{HTML}{9588A5}\definecolor{VPF}{HTML}{ECE8F0}
\definecolor{VI}{HTML}{3B4252}\definecolor{VM}{HTML}{6B7280}\definecolor{VL}{HTML}{4C4D4F}
\begin{tikzpicture}[x=1mm,y=1mm,>=Stealth,font=\figurefont\fontsize{8}{10}\selectfont,
lab/.style={text=VM,font=\figurefont\fontsize{7.5}{9}\selectfont}]
\def\L{80}
\draw[VL,line width=.6pt] (0,0) -- (\L,0);
\foreach \x/\t in {0/0,1/1} { \draw[VL,line width=.6pt] ({\L*\x},0) -- ({\L*\x},1.2); \node[below,text=VI] at ({\L*\x},-.3) {$\t$}; }
% rows: centre, k, length
\foreach \c/\k/\len/\lab in {0.5/1/0.1/{1/2}, 0.333333/2/0.05/{1/3}, 0.666667/3/0.025/{2/3}, 0.25/4/0.0125/{1/4}, 0.75/5/0.00625/{3/4}} {
  \pgfmathsetmacro{\y}{3+4.2*\k}
  \fill[VBF] ({\L*(\c-\len/2)},{\y-1}) rectangle ({\L*(\c+\len/2)},{\y+1});
  \draw[VB,line width=.6pt] ({\L*(\c-\len/2)},{\y-1}) rectangle ({\L*(\c+\len/2)},{\y+1});
  \draw[VM,densely dotted,line width=.4pt] ({\L*\c},0) -- ({\L*\c},{\y-1});
  \fill[VR] ({\L*\c},0) circle (.75);
  \node[anchor=west,text=VI] at ({\L+4},\y) {$q_{\k}=\lab$};
}
\node[anchor=west,text=VM] at ({\L+4},{3+4.2*6}) {$\vdots$};
\node[anchor=west,lab] at ({\L+4},0) {长度合计不超过$\varepsilon$};
\end{tikzpicture}
